% ============================================================================
%  CAPÍTULO II: DIAGNÓSTICO DEL ORIGEN Y HARDENING DEL DESTINO
% ============================================================================

\section{Arquitectura del sistema de origen}

El sistema de origen es una instalación MySQL 5.7.44 sobre Linux, ejecutada
en el contenedor \texttt{coop\_horizonte\_legacy}. La base se denomina
\texttt{coop\_horizonte\_legacy}, usa el motor InnoDB y la collation
\texttt{utf8\_general\_ci}. La aplicación legacy se conecta con el usuario
\texttt{app\_legacy\_user}, creado con \texttt{GRANT ALL PRIVILEGES} sobre
toda la base. No existe una cuenta separada para auditoría, no se exige TLS
y no hay un esquema de respaldo cifrado dentro del flujo operativo.

La Figura~\ref{fig:arquitectura-origen} representa la arquitectura
encontrada. La aplicación y el motor comparten la red del contenedor; el
tráfico no está cifrado y el usuario de aplicación puede ejecutar DDL, DML y
DCL. El puerto 3306 se publica únicamente en \texttt{127.0.0.1}, lo que
reduce la exposición desde la red del anfitrión, pero no corrige los
problemas internos de privilegios ni la falta de cifrado.

\figuraAPA{fig:arquitectura-origen}{Arquitectura del sistema de origen MySQL 5.7}{%
  \begin{tikzpicture}[
      nodo/.style={draw, rounded corners, minimum width=3.2cm,
                   minimum height=1.0cm, align=center, font=\small},
      riesgo/.style={draw, dashed, rounded corners, minimum width=3.2cm,
                     minimum height=0.8cm, align=center, font=\scriptsize,
                     text=black!70},
      linea/.style={-Stealth, thick}]
    \node[nodo] (app) {Aplicación\\legacy};
    \node[nodo, right=2.4cm of app] (db) {MySQL 5.7.44\\\texttt{coop\_horizonte\_legacy}};
    \node[riesgo, below=0.9cm of app] (r1) {Sin TLS};
    \node[riesgo, below=0.9cm of db] (r2) {\texttt{GRANT ALL}};
    \draw[linea] (app) -- node[above, font=\scriptsize]{TCP/3306} (db);
    \draw[dashed, red!60] (r1) -- (app);
    \draw[dashed, red!60] (r2) -- (db);
  \end{tikzpicture}%
}{Elaboración propia con datos de \texttt{evidence/01\_inventario\_origen.txt}.}

\capturaAPA{fig:cap-origen-version}{Versión y base de datos del sistema de origen}%
  {captura-origen-version.png}%
  {Salida de \texttt{SELECT VERSION();} y \texttt{SHOW DATABASES;} en MySQL 5.7.44.}

El origen contiene 8 tablas de negocio, 2 vistas, 19 entradas de índice y un
solo usuario de aplicación. La Tabla~\ref{tab:inventario-origen} resume el
inventario. El tamaño total reportado por \texttt{information\_schema} es de
0.28 MB, un volumen pequeño que permite una ventana de migración corta y un
enfoque de corte controlado (\emph{big bang}) en lugar de replicación continua.

\begin{table}[H]
  \centering
  \caption{Inventario de objetos del sistema de origen}
  \label{tab:inventario-origen}
  \contenidoConFuente{%
    \begin{tablaAPA}
      \begin{tabular}{@{}L{4.6cm}C{2.0cm}L{7.9cm}@{}}
        \toprule
        \textbf{Tipo de objeto} & \textbf{Cantidad} & \textbf{Detalle} \\
        \midrule
        Tablas base & 8 & \texttt{sucursales}, \texttt{socios}, \texttt{cuentas\_ahorro}, \texttt{tipos\_credito}, \texttt{solicitudes\_credito}, \texttt{prestamos}, \texttt{cuotas\_amortizacion}, \texttt{transacciones}. \\
        Vistas & 2 & \texttt{vw\_cartera\_mora} y \texttt{vw\_resumen\_saldos\_socio}. \\
        Entradas de índice & 19 & 8 claves primarias, 3 \texttt{UNIQUE} y 8 índices no únicos (uno compuesto). \\
        Secuencias & 0 & MySQL 5.7 no define objetos \texttt{SEQUENCE}; 8 columnas \texttt{AUTO\_INCREMENT} pasan a identidades en PostgreSQL. \\
        Usuarios de aplicación & 1 & \texttt{app\_legacy\_user} con \texttt{GRANT ALL}. \\
        Claves foráneas & 0 & La integridad referencial no está declarada en el motor. \\
        Tamaño de la base & 0.28 MB & Medido con \texttt{data\_length + index\_length}. \\
        \bottomrule
      \end{tabular}
    \end{tablaAPA}%
  }{Elaboración propia con datos de \texttt{evidence/01\_inventario\_origen.txt}.}
\end{table}

La ausencia de objetos \texttt{SEQUENCE} no implica que el destino carezca de
generadores de identificadores: las 8 columnas \texttt{AUTO\_INCREMENT} se
convierten en identidades \texttt{BIGINT GENERATED ALWAYS AS IDENTITY} y sus
valores se reajustan con \texttt{setval} al final de la carga, sin inventar
objetos nuevos ni perder los identificadores originales. {\hyphenation{OVERRIDING}%
La documentación oficial de PostgreSQL 16 describe la identidad así: «In an
INSERT command, if ALWAYS is selected, a user-specified value is only accepted
if the INSERT statement specifies OVERRIDING SYSTEM VALUE»
\citep[apartado «GENERATED ALWAYS AS IDENTITY»]{postgresql2024}.}

\capturaAPA{fig:cap-inventario}{Inventario de objetos del sistema de origen}%
  {captura-inventario-objetos.png}%
  {Salida de \texttt{scripts/01\_audit\_origin.sh} con tablas, vistas, índices, usuarios y privilegios.}

La salida completa del inventario se presenta en el Anexo~C
(Figuras~\ref{fig:anexo-inventario-1} y~\ref{fig:anexo-inventario-2}).

\subsection{Diagrama de clases de la base de datos origen}

La Figura~\ref{fig:diagrama-clases-origen} presenta el diagrama de clases de
la base de datos origen. Cada clase representa una tabla de
\texttt{coop\_horizonte\_legacy}, con sus atributos y tipos de datos de
MySQL 5.7. Las asociaciones reflejan las relaciones lógicas del negocio: un
socio tiene muchas cuentas de ahorro y muchas solicitudes de crédito; una
solicitud aprobada genera un préstamo; un préstamo tiene muchas cuotas de
amortización; y las transacciones referencian un socio, una cuenta o un
préstamo según el tipo de movimiento. Estas relaciones son lógicas, inferidas
del modelo de datos y del código de la aplicación: como se documenta en la
Sección~\ref{sec:riesgos}, el origen no declara claves foráneas en el motor,
y solo se formalizan como restricciones \texttt{FOREIGN KEY} en el destino.

\noindent\begin{minipage}{\linewidth}
\capturaAPA{fig:diagrama-clases-origen}{Diagrama de clases de la base de datos origen}%
  {diagrama-clases-origen.png}%
  {Elaboración propia con base en \texttt{src/origin/01\_schema\_origin.sql}.}
\end{minipage}

\section{Inventario de datos y volumen}

El conjunto de datos fue generado de forma determinista con una semilla fija
(\texttt{seed=42}) para que cualquier integrante obtenga exactamente los
mismos valores. La Tabla~\ref{tab:conteos-origen} presenta los conteos
exactos. Destacan los 55 socios, las 66 cuentas de ahorro (algunos socios
tienen dos cuentas), las 50 solicitudes de crédito, los 32 préstamos
aprobados, las 516 cuotas de amortización y las 740 transacciones. El total
es de 1\,469 filas.

\begin{table}[H]
  \centering
  \caption{Conteo de filas en el sistema de origen}
  \label{tab:conteos-origen}
  \contenidoConFuente{%
    \begin{tablaAPA}
      \begin{tabular}{@{}L{6.2cm}R{3.0cm}L{5.0cm}@{}}
        \toprule
        \textbf{Tabla} & \textbf{Filas} & \textbf{Propósito} \\
        \midrule
        \texttt{sucursales} & 5 & Catálogo de agencias. \\
        \texttt{socios} & 55 & Datos personales y estado del socio. \\
        \texttt{cuentas\_ahorro} & 66 & Cuentas y saldos iniciales. \\
        \texttt{tipos\_credito} & 5 & Productos de crédito. \\
        \texttt{solicitudes\_credito} & 50 & Solicitudes y su estado. \\
        \texttt{prestamos} & 32 & Créditos desembolsados. \\
        \texttt{cuotas\_amortizacion} & 516 & Plan de pagos. \\
        \texttt{transacciones} & 740 & Depósitos, retiros, pagos y desembolsos. \\
        \midrule
        \textbf{Total} & \textbf{1\,469} & 8 tablas de negocio. \\
        \bottomrule
      \end{tabular}
    \end{tablaAPA}%
  }{Elaboración propia con datos de \texttt{evidence/01\_inventario\_origen.txt} y \texttt{evidence/02\_ejecucion\_migracion.txt}.}
\end{table}

La distribución de filas no es arbitraria: las transacciones se derivan de
los desembolsos, los pagos de cuotas y los movimientos de las cuentas; las
cuotas se generan a partir de cada préstamo; y los préstamos solo existen
para solicitudes aprobadas. Esa coherencia es la que permite verificar la
integridad referencial después de la migración. El \emph{script}
\texttt{src/origin/generate\_seed.py} documenta las reglas de generación y
puede regenerar el archivo \texttt{02\_seed\_origin.sql} en cualquier momento.

\section{Riesgos de seguridad detectados}
\label{sec:riesgos}

\subsection{Privilegios excesivos del usuario de aplicación}

El hallazgo más grave es que \texttt{app\_legacy\_user} tiene
\texttt{GRANT ALL PRIVILEGES} sobre \texttt{coop\_horizonte\_legacy.*}. En
MySQL, \texttt{ALL} incluye \texttt{SELECT}, \texttt{INSERT}, \texttt{UPDATE},
\texttt{DELETE}, \texttt{CREATE}, \texttt{DROP}, \texttt{ALTER},
\texttt{INDEX}, \texttt{REFERENCES}, \texttt{CREATE TEMPORARY TABLES},
\texttt{LOCK TABLES}, \texttt{EXECUTE}, \texttt{CREATE VIEW},
\texttt{SHOW VIEW}, \texttt{CREATE ROUTINE}, \texttt{ALTER ROUTINE},
\texttt{EVENT} y \texttt{TRIGGER}. La aplicación no necesita DDL ni
\texttt{DROP}; sin embargo, si es comprometida, el atacante puede borrar
tablas, alterar el esquema o crear disparadores maliciosos. La mitigación es
separar la operación de la administración y otorgar solo DML sobre las tablas
necesarias, que es exactamente lo que hace el rol \texttt{app\_core\_rw} en
el destino.

\capturaAPA{fig:cap-privilegios}{Privilegios del usuario de aplicación en el origen}%
  {captura-privilegios-app.png}%
  {Salida de \texttt{SHOW GRANTS FOR 'app\_legacy\_user'@'\%';} con \texttt{GRANT ALL PRIVILEGES}.}

\subsection{Transporte sin cifrado}

MySQL 5.7 puede habilitar TLS, pero la instalación analizada no lo exige. La
aplicación se conecta sin \texttt{REQUIRE SSL}, de modo que las credenciales
y los datos pueden ser interceptados en la red interna. En el destino se
exige \texttt{hostssl} para todas las conexiones TCP y se rechaza cualquier
intento sin cifrado, como se demuestra en la
Sección~\ref{sec:evidencias-hardening}.

\subsection{Ausencia de integridad referencial}

El esquema no declara claves foráneas. La aplicación es la única responsable
de que \texttt{cuentas\_ahorro.socio\_id} apunte a un socio existente, que
\texttt{prestamos.solicitud\_id} corresponda a una solicitud aprobada y que
\texttt{cuotas\_amortizacion.prestamo\_id} pertenezca a un préstamo. Cualquier
error de código, carga manual o borrado accidental puede producir huérfanos.
En el destino se añadieron 9 claves foráneas y se verificó que no existieran
huérfanos después de la carga.

\subsection{Auditoría insuficiente}

No hay registro sistemático de conexiones, desconexiones ni cambios de
estructura. Si ocurriera un acceso no autorizado o un cambio de esquema, la
cooperativa no tendría evidencia para reconstruir los hechos. El destino
habilita \texttt{log\_connections}, \texttt{log\_disconnections} y
\texttt{log\_statement = 'ddl'}, y además separa los roles para que la cuenta
de auditoría no pueda modificar datos.

\subsection{Respaldo sin cifrado y sin prueba de restauración}

Un respaldo en texto claro es tan sensible como la base original. Si el
archivo se copia a un medio externo, se comparte por error o se almacena en
un servidor comprometido, los datos personales y financieros quedan
expuestos. El proyecto genera el respaldo del destino con
\texttt{pg\_dump -Fc} y lo cifra con AES-256-CBC y PBKDF2; además, descifra y
restaura el respaldo en una base temporal para comprobar que es utilizable.
Un respaldo que nunca se restaura no es un respaldo: es una esperanza.

\subsection{Fin de vida del motor}

MySQL 5.7 dejó de recibir parches oficiales en octubre de 2023. La
cooperativa no puede depender de un motor sin soporte para custodiar datos
financieros. PostgreSQL 16 tiene soporte activo, una comunidad amplia y un
ciclo de versiones predecible \citep{postgresql2024}. La migración reduce el
riesgo de vulnerabilidades no corregidas y facilita la aplicación de futuras
actualizaciones.

\section{Matriz de incompatibilidades y reglas de conversión}

La migración heterogénea no consiste en copiar valores: exige interpretar la
semántica de cada tipo. La Tabla~\ref{tab:incompatibilidades} presenta las
incompatibilidades reales encontradas, su frecuencia y la regla aplicada en el
destino. Las decisiones se justifican con la documentación oficial de
PostgreSQL \citep{postgresql2024} y con los principios de diseño relacional
\citep{elmasri2016,silberschatz2020}.

\begin{table}[H]
  \centering
  \caption{Matriz de incompatibilidades MySQL 5.7--PostgreSQL 16}
  \label{tab:incompatibilidades}
  \contenidoConFuente{%
    \begin{tablaAPA}
      \begin{tabular}{@{}L{2.7cm}L{3.3cm}L{3.6cm}L{5.2cm}@{}}
        \toprule
        \textbf{MySQL 5.7} & \textbf{Columnas reales} & \textbf{PostgreSQL 16} & \textbf{Regla de conversión} \\
        \midrule
        \texttt{TINYINT(1)} & 4 columnas: \texttt{estatus} en \texttt{sucursales}, \texttt{socios} y \texttt{cuentas\_ahorro}; \texttt{pagado\_total} en \texttt{prestamos}. & \texttt{BOOLEAN} & Convertir 0 a \texttt{FALSE} y cualquier valor distinto de 0 a \texttt{TRUE}. Validar que no existan valores intermedios. \\
        \texttt{ENUM} & 7 columnas: estados de socios, tipos de crédito, solicitudes, préstamos, cuotas y transacciones. & \texttt{VARCHAR(n)} + \texttt{CHECK} & Reemplazar el tipo por texto con una restricción \texttt{CHECK ... IN (...)} que conserve el dominio. \\
        \texttt{DATETIME} & 10 columnas: fechas de creación, apertura, cierre, solicitud, aprobación, desembolso, pago y transacción. & \texttt{TIMESTAMPTZ} & Interpretar el valor original en \texttt{America/La\_Paz} y almacenarlo como instante con zona horaria. \\
        \texttt{INT UNSIGNED AUTO\_INCREMENT} & 8 claves primarias. & \texttt{BIGINT GENERATED ALWAYS AS IDENTITY} & Crear identidad, insertar los IDs originales con \texttt{OVERRIDING SYSTEM VALUE} y reajustar la secuencia con \texttt{setval}. \\
        \texttt{utf8\_general\_ci} & Base y tablas. & UTF-8 nativo & Usar la codificación UTF-8 del clúster; no se requiere collation por columna. \\
        \texttt{UNIQUE KEY} & 3 columnas únicas: \texttt{numero\_identificacion}, \texttt{num\_cuenta}, \texttt{solicitud\_id}. & \texttt{UNIQUE} & Conservar la restricción y crear el índice único asociado. \\
        Sin \texttt{FOREIGN KEY} & 0 restricciones en el origen. & 9 \texttt{FOREIGN KEY} & Declarar las relaciones y validar que no existan huérfanos antes de la carga. \\
        \texttt{DECIMAL(12,2)} & Montos, saldos, tasas y cuotas. & \texttt{NUMERIC(12,2)} & Mantener precisión y escala; añadir \texttt{CHECK} para montos no negativos. \\
        \texttt{TEXT} & Descripciones, observaciones y direcciones. & \texttt{TEXT} & Conservar sin cambios; PostgreSQL no requiere prefijos de longitud. \\
        \bottomrule
      \end{tabular}
    \end{tablaAPA}%
  }{Elaboración propia con datos de \texttt{evidence/01\_inventario\_origen.txt} y \texttt{src/target/01\_schema\_target.sql}.}
\end{table}

\subsection{Incompatibilidades de sintaxis y adaptación}

Además de los tipos, la sintaxis de ambos motores difiere. La adaptación no
se resolvió con una conversión automática: se revisaron las sentencias del
origen y se aplicaron las equivalencias de la
Tabla~\ref{tab:incompatibilidades-sintaxis}.

\begin{table}[H]
  \centering
  \caption{Incompatibilidades de sintaxis MySQL 5.7--PostgreSQL 16}
  \label{tab:incompatibilidades-sintaxis}
  \contenidoConFuente{%
    \begin{tablaAPA}
      \begin{tabular}{@{}L{4.0cm}L{4.3cm}L{6.5cm}@{}}
        \toprule
        \textbf{MySQL 5.7} & \textbf{PostgreSQL 16} & \textbf{Decisión} \\
        \midrule
        \code{ENGINE=InnoDB}, \code{CHARSET}, \code{COLLATE} & Sintaxis nativa & Se omiten; se usa la codificación UTF-8 del clúster. \\
        \code{AUTO\_INCREMENT} & \code{GENERATED ALWAYS AS IDENTITY} & Se conservan los IDs y se reajusta con \texttt{setval}. \\
        \code{GRANT ALL ON db.*} con comillas invertidas & Roles y \texttt{GRANT} sobre el esquema & Se separan \texttt{admin\_dba}, \texttt{app\_core\_rw} y \texttt{auditor\_ro}. \\
        \code{SHOW GRANTS} & \texttt{pg\_roles} y \texttt{has\_table\_privilege} & Los privilegios se auditan por catálogo. \\
        \code{IFNULL} & \code{COALESCE} & Se adopta la función estándar SQL. \\
        \code{GROUP\_CONCAT} & \code{string\_agg} & Se usa en las vistas de reporte. \\
        \code{ON DUPLICATE KEY UPDATE} & \code{INSERT ... ON CONFLICT} & Disponible para cargas tipo \emph{upsert}. \\
        Comentarios con \texttt{\#} & Comentarios con \texttt{-{}-} & Se convierten al adaptar los scripts. \\
        Comillas invertidas para identificadores & Comillas dobles o sin delimitar & Se prefiere SQL estándar. \\
        \bottomrule
      \end{tabular}
    \end{tablaAPA}%
  }{Elaboración propia con base en \texttt{src/origin/01\_schema\_origin.sql} y \texttt{src/target/01\_schema\_target.sql}.}
\end{table}

Estas diferencias explican por qué el DDL y el ETL se escribieron
explícitamente en lugar de generarse con una herramienta genérica.

\subsection{Booleanos: de \texttt{TINYINT(1)} a \texttt{BOOLEAN}}

MySQL no tiene un tipo booleano nativo: \texttt{BOOLEAN} es un alias de
\texttt{TINYINT(1)} y acepta cualquier valor de 8 bits. PostgreSQL tiene
\texttt{BOOLEAN} con valores \texttt{TRUE}, \texttt{FALSE} y \texttt{NULL}.
La conversión aplicada en el ETL es explícita: \texttt{bool(int(valor))}. Si
el origen contuviera un 2 o un -1, el ETL lo interpretaría como verdadero; por
eso se recomienda una validación previa en producción. En el conjunto de
datos analizado todos los valores son 0 o 1.

\subsection{Estados: de \texttt{ENUM} a \texttt{VARCHAR} con \texttt{CHECK}}

\texttt{ENUM} es una enumeración almacenada como entero interno. Su principal
ventaja es la validación en el motor; su principal desventaja es que añadir un
valor requiere \texttt{ALTER TABLE} y que el orden interno puede cambiar si se
reordena la lista. En PostgreSQL se optó por \texttt{VARCHAR(10)} con una
restricción \texttt{CHECK}. Así se conserva la validación, se facilita la
lectura y se evita depender de un tipo propietario. Las restricciones se
nombraron de forma descriptiva para que los errores sean comprensibles.

\subsection{Fechas: de \texttt{DATETIME} a \texttt{TIMESTAMPTZ}}

\texttt{DATETIME} almacena una fecha y hora sin zona. Eso significa que el
valor \texttt{2026-11-20 08:09:40} no indica si ocurrió en La Paz, en UTC o en
otra zona. Para una entidad financiera, la ambigüedad es inaceptable: un pago
registrado a las 23:30 puede pertenecer a un día contable distinto según la
zona. La migración interpreta cada valor en \texttt{America/La\_Paz} y lo
almacena como \texttt{TIMESTAMPTZ}. El ETL usa \texttt{ZoneInfo} de Python
para asignar la zona antes de insertar, y el servidor configura
\texttt{timezone = 'America/La\_Paz'} para que las consultas y los reportes
se muestren en la hora local.

\subsection{Identificadores: de \texttt{AUTO\_INCREMENT} a identidad}

MySQL usa \texttt{AUTO\_INCREMENT} con enteros sin signo. PostgreSQL 16 usa
identidades (\texttt{GENERATED ALWAYS AS IDENTITY}), que son estándar SQL y
permiten controlar el ciclo de vida de la secuencia. El ETL inserta los IDs
originales con \texttt{OVERRIDING SYSTEM VALUE} para no romper las referencias
y, al final, ejecuta \texttt{setval} sobre cada secuencia. Sin ese ajuste, el
primer \texttt{INSERT} nuevo intentaría usar un ID ya ocupado y fallaría.

\subsection{Collation y codificación}

\texttt{utf8\_general\_ci} es la collation de \texttt{utf8mb3}, una
codificación de tres bytes que no puede representar todos los caracteres
Unicode (por ejemplo, emojis o ciertos símbolos). PostgreSQL usa UTF-8 nativo
y una collation del clúster. La migración no requiere conversión de datos
porque los valores del origen son ASCII y español; sin embargo, el cambio
elimina la limitación y prepara la base para nombres y direcciones con
caracteres adicionales.

\subsection{Integridad referencial}

La ausencia de claves foráneas en el origen obligó a validar la integridad
antes de declararlas en el destino. El ETL carga en orden jerárquico
(\texttt{sucursales}, \texttt{tipos\_credito}, \texttt{socios},
\texttt{cuentas\_ahorro}, \texttt{solicitudes\_credito}, \texttt{prestamos},
\texttt{cuotas\_amortizacion}, \texttt{transacciones}) y luego las pruebas
consultan cada relación con \texttt{LEFT JOIN} para contar huérfanos. El
resultado fue cero en las 9 relaciones, como se muestra en el
Capítulo~IV.
